% !TeX root = booklet-02.tex
\chapter*{تمرین‌های اثر آستانه و مرور روابط}
\addcontentsline{toc}{chapter}{تمرین‌های اثر آستانه و مرور روابط}
در این فصل، تمرین‌های \secnum{۲.۳} و \secnum{۲.۷} تا \secnum{۲.۹} بر اساس صورت اصلی کتاب و با
استفاده از روابط بخش \ref{sec:threshold-effects} حل می‌شوند.
حل تمرین \secnum{۲.۳}، پیش‌نیاز مدار سری تمرین \secnum{۲.۲} را نیز در حد لازم بازسازی
می‌کند تا خواندن استدلال به فایل دیگری وابسته نباشد.

\section*{تمرین \secnum{۲.۳}، مقایسه جریان دو ترانزیستور سری}
\paragraph*{صورت مسئله و مدار}
در تمرین \secnum{۲.۲}، جریان دو ترانزیستور یکسان سری با جریان یک ترانزیستور
به طول دو برابر مقایسه شده است. تمرین \secnum{۲.۳} می‌پرسد اگر اثر بدنه
لحاظ شود، جریان ساختار سری چه نسبتی با جریان ترانزیستور منفرد دارد
\cite[تمرین‌های \secnum{۲.۲} و \secnum{۲.۳}، صفحه ۹۵]{weste}.
مدارهای لازم در شکل \ref{fig:body-series-comparison} بازترسیم شده‌اند.
همه بدنه‌ها را به زمین وصل می‌کنیم و مدل \lr{Shockley} با ضریب
\(\beta=\mu C_{\mathrm{ox}}W/L\) برای هر ترانزیستور طول \(L\) به کار
می‌بریم. آستانه در بایاس صفر \(V_{t0}\) و ضریب اثر بدنه \(\gamma>0\)
است. اثرهای \lr{DIBL} و \lr{CLM} در این مقایسه ثابت و صفر فرض می‌شوند.

\begin{figure}[!htbp]
\centering
\begin{latin}
\PersianMathDigits
\begin{circuitikz}[scale=0.9,transform shape]
 \draw (1.8,2) node[nmos] (single) {};
 \draw (single.S)--(1.8,0.4) node[ground] {};
 \draw (single.D)--(1.8,4.3) node[ocirc] {} node[above] {\(D\)};
 \draw (single.G)--(0,2) node[ocirc] {} node[left] {\(V_{DD}\)};
 \node[right] at (2.2,1.2) {\(W/(2L)\)};
 \draw[->] (2.9,3.8)--(2.9,3.1) node[midway,right] {\(I_1\)};
 \draw (7.2,1.3) node[nmos] (bottom) {};
 \draw (7.2,3.3) node[nmos] (top) {};
 \draw (bottom.S)--(7.2,0.4) node[ground] {};
 \draw (bottom.D)--(top.S);
 \draw (top.D)--(7.2,4.3) node[ocirc] {} node[above] {\(D\)};
 \draw (bottom.G)--(5.4,1.3)--(5.4,3.3)--(top.G);
 \draw (5.4,2.3)--(4.6,2.3) node[ocirc] {} node[left] {\(V_{DD}\)};
 \draw (7.2,2.3) node[circ] {}--(8.3,2.3) node[right] {\(x\)};
 \node[right] at (7.6,1.3) {\(W/L\)};
 \node[right] at (7.6,3.3) {\(W/L\)};
 \draw[->] (8.5,4.2)--(8.5,3.6) node[midway,right] {\(I_2\)};
\end{circuitikz}
\end{latin}
\caption{بازترسیم آموزشی شکل \secnum{۲.۳۲} کتاب برای تمرین \secnum{۲.۳}.
\(D\) ولتاژ درین بالایی و \(x\) ولتاژ گره میانی نسبت به زمین است.
همه بدنه‌ها به زمین متصل‌اند \cite[صفحه ۹۵]{weste}.}
\label{fig:body-series-comparison}
\end{figure}

\paragraph*{بازسازی نتیجه مرجع با آستانه ثابت}
برای اختصار، \(A=V_{DD}-V_{t0}>0\) می‌گذاریم. ابتدا دامنه اصلی
\(0<D<A\) را بررسی می‌کنیم. ضریب ترانزیستور منفرد طول \(2L\)،
\(\beta/2\) است. بنابراین جریان آن برابر است با
\begin{equation}
 I_1=\frac{\beta}{2}\left(AD-\frac{D^2}{2}\right).
 \label{eq:body-stack-single}
\end{equation}
برای ساختار سری، \(V_{GS}\) و \(V_{DS}\) ترانزیستور پایین به‌ترتیب
\(V_{DD}\) و \(x\) هستند. برای ترانزیستور بالا، این دو ولتاژ
\(V_{DD}-x\) و \(D-x\) هستند. با آستانه ثابت، جریان‌های دو ترانزیستور
از رابطه \lr{Linear} به دست می‌آیند
\begin{align}
 f(x)&=\beta\left(Ax-\frac{x^2}{2}\right),\notag\\
 g_0(x)&=\beta\left[(A-x)(D-x)-\frac{(D-x)^2}{2}\right]\notag\\
 &=\beta\left(AD-\frac{D^2}{2}-Ax+\frac{x^2}{2}\right).
 \label{eq:body-stack-reference-currents}
\end{align}
در حالت ماندگار، جریان واردشده و خارج‌شده از گره میانی برابرند.
این همان \lr{Kirchhoff's Current Law (KCL)} است، پس \(f(x)=g_0(x)\).
جمع دو رابطه نیز \(f(x)+g_0(x)=\beta(AD-D^2/2)\) می‌دهد. بنابراین
جریان هر ترانزیستور نصف این مجموع و برابر \(I_1\) است.

با برابرکردن دو جریان و جمع‌کردن جمله‌های هم‌نوع، داریم
\begin{align}
 Ax-\frac{x^2}{2}&=AD-\frac{D^2}{2}-Ax+\frac{x^2}{2},\notag\\
 x^2-2Ax+AD-\frac{D^2}{2}&=0,\notag\\
 x&=A\pm\sqrt{A^2-AD+\frac{D^2}{2}}.
 \label{eq:body-stack-reference-roots}
\end{align}
ولتاژ فیزیکی گره میانی باید در بازه \(0<x<D\) قرار گیرد.
ریشه با علامت مثبت از \(A>D\) بزرگ‌تر می‌شود. ریشه پذیرفته‌شده چنین است
\begin{equation}
 x_0=A-\sqrt{A^2-AD+\frac{D^2}{2}},\qquad 0<x_0<D.
 \label{eq:body-stack-reference-root}
\end{equation}
برای بررسی بازه، عبارت زیر رادیکال بین \((A-D)^2\) و \(A^2\) قرار
دارد، زیرا \(0<D<A\). با ریشه‌گرفتن،
\(A-D<\sqrt{A^2-AD+D^2/2}<A\) و در نتیجه \(0<x_0<D\) حاصل می‌شود.
پس هر دو ترانزیستور در حل مرجع روشن و در ناحیه \lr{Linear} هستند.

\paragraph*{افزودن اثر بدنه و تغییر نقطه کار}
در ترانزیستور منفرد و ترانزیستور پایین، سورس و بدنه در زمین‌اند.
بنابراین آستانه آن‌ها همان \(V_{t0}\) می‌ماند. در ترانزیستور بالا،
\(V_{SB}=x>0\) است و افزایش آستانه برابر است با
\begin{equation}
 \delta(x)=\gamma\left(\sqrt{\phi_s+x}-\sqrt{\phi_s}\right)>0.
 \label{eq:body-stack-shift}
\end{equation}
در یک \(x\) ثابت، اضافه‌ولتاژ ترانزیستور بالا از \(A-x\) به
\(A-x-\delta(x)\) کاهش می‌یابد. جریان آن را با لحاظ اثر بدنه
\(g_{\mathrm{B}}(x)\) می‌نامیم. رابطه جریان در ناحیه \lr{Linear}
یا \lr{Saturation} بر اساس همین اضافه‌ولتاژ جدید انتخاب می‌شود.

در نقطه مرجع \(x_0\)، ترانزیستور پایین همان جریان \(I_1\) را دارد،
اما جریان ترانزیستور بالا به دلیل آستانه بیشتر کاهش یافته است.
پس برابری جریان‌ها در \(x_0\) برقرار نمی‌ماند و ولتاژ گره میانی باید
تغییر کند. در بازه \(0<x<D<A\)، جریان پایین با افزایش \(x\) زیاد
می‌شود، زیرا شیب آن \(f'(x)=\beta(A-x)>0\) است.
جریان بالای مرجع با افزایش \(x\) کم می‌شود، زیرا از
\eqref{eq:body-stack-reference-currents}،
\(g_0'(x)=-\beta(A-x)<0\) به دست می‌آید.

برای هر \(x\ge x_0\)، جریان بالای جدید از جریان مرجع در همان ولتاژ
کمتر است و جریان مرجع نیز حداکثر \(g_0(x_0)=I_1\) است.
در همین بازه، جریان پایین دست‌کم \(f(x_0)=I_1\) است.
بنابراین نقطه برابری جریان جدید در \(x_{\mathrm{B}}<x_0\) قرار می‌گیرد.
در \(x=0\)، جریان پایین صفر و جریان بالا مثبت است، پس تغییر از این
نقطه به \(x_0\) یک نقطه برابری در میان آن‌ها ایجاد می‌کند.
نتیجه چنین است
\begin{equation}
 \boxed{I_{2,\mathrm{B}}=f(x_{\mathrm{B}})<f(x_0)=I_1}.
 \label{eq:body-stack-current-result}
\end{equation}
این استدلال، جابه‌جایی گره میانی را نیز وارد مقایسه می‌کند.
اگر ترانزیستور بالا در نقطه جدید اشباع شود، کاهش اضافه‌ولتاژ همچنان
جریان آن را در \(x\) ثابت کاهش می‌دهد و مقایسه برقرار است.

در دامنه \(D\ge A\)، ترانزیستور بالای روشن در اشباع است و جریان آن
\(\beta[A-x-\delta(x)]^2/2\) می‌شود. با آستانه ثابت، جریان مرجع
ساختار سری و ترانزیستور منفرد هر دو \(\beta A^2/4\) هستند.
افزایش آستانه بالایی باز هم نقطه کار را به \(x\) کوچک‌تر و جریان
کمتر جابه‌جا می‌کند. در حد \(D=0\)، دو جریان صفر و برابرند.
در حد \(\gamma=0\)، نتیجه برابری تمرین \secnum{۲.۲} بازیابی می‌شود.

\paragraph*{مثال عددی آموزشی برای جابه‌جایی گره میانی}
فرض کنید \(A=\qty{1}{\volt}\)، \(D=\qty{0.5}{\volt}\)،
\(\beta=\qty{200}{\micro\ampere\per\volt\squared}\)،
\(\gamma=\qty{0.5}{\volt\tothe{0.5}}\) و
\(\phi_s=\qty{0.8}{\volt}\) باشند. رابطه
\eqref{eq:body-stack-reference-root} می‌دهد
\(x_0\approx\qty{0.20943}{\volt}\). جریان مرجع نیز برابر است با
\[
 I_1=(200/2)[1(0.5)-0.5^2/2]=\qty{37.5}{\micro\ampere}.
\]

برای محاسبه با اثر بدنه، مقدار آزمایشی \(x\) را در
\(f(x)\) و جریان بالایی قرار می‌دهیم و با کوچک‌کردن بازه اختلاف
جریان‌ها به نقطه برابری می‌رسیم. نتیجه عددی چنین است
\begin{equation}
 x_{\mathrm{B}}\approx\qty{0.19948}{\volt},\qquad
 I_{2,\mathrm{B}}\approx\qty{35.92}{\micro\ampere}.
 \label{eq:body-stack-numeric}
\end{equation}
در این نقطه، \(\delta(x_{\mathrm{B}})\approx\qty{0.05266}{\volt}\)،
اضافه‌ولتاژ بالا حدود \(\qty{0.74786}{\volt}\) و ولتاژ درین و سورس
آن حدود \(\qty{0.30052}{\volt}\) است. هر دو ترانزیستور در
\lr{Linear} قرار دارند. جریان حدود \num{4.22} درصد کاهش یافته است.
این مثال، محاسبه عددی معادله مدل است.

\section*{تمرین \secnum{۲.۷}، محاسبه آستانه در بایاس سورس بزرگ}
\paragraph*{صورت مسئله و داده‌ها}
ترانزیستور \lr{nMOS} در فرایند \(\qty{0.6}{\micro\metre}\) دارای
ضخامت اکسید صد آنگستروم،
\(N_A=\qty{2e17}{\per\centi\metre\cubed}\) و آستانه مرجع
\(V_{t0}=\qty{0.7}{\volt}\) است. بدنه به زمین متصل است. خواسته،
تغییر آستانه در دمای اتاق هنگام افزایش ولتاژ سورس از صفر به
\(\qty{4}{\volt}\) است \cite[تمرین \secnum{۲.۷}، صفحه ۹۶]{weste}.
از \(V_{SB}=V_S-V_B\)، بایاس جدید برابر \(\qty{4}{\volt}\) است.
ثابت‌های تقریبی و قرارداد دما را مانند مثال
\eqref{eq:body-book-example-phi} انتخاب می‌کنیم.

\paragraph*{تبدیل ضخامت و محاسبه پتانسیل سطحی}
هر آنگستروم برابر \(\qty{e-10}{\metre}\) است. پس ضخامت برابر
\(\qty{10}{\nano\metre}=\qty{e-6}{\centi\metre}\) می‌شود.
با \(v_T\approx\qty{0.026}{\volt}\) و
\(n_i=\qty{1.45e10}{\per\centi\metre\cubed}\)، داریم
\begin{align}
 \frac{N_A}{n_i}&=\frac{2\times10^{17}}{1.45\times10^{10}}
 \approx1.3793\times10^7,\notag\\
 \ln\!\left(\frac{N_A}{n_i}\right)&\approx16.4397,\notag\\
 \phi_s&\approx2(0.026)(16.4397)\,\unit{\volt}
 \approx\qty{0.85486}{\volt}.
 \label{eq:body-ex7-phi}
\end{align}
ظرفیت اکسید بر واحد سطح از تقسیم \lr{Permittivity} اکسید بر ضخامت
به دست می‌آید
\begin{equation}
 C_{\mathrm{ox}}=\frac{3.9(8.85\times10^{-14})}{10^{-6}}
 =\qty{3.4515e-7}{\farad\per\centi\metre\squared}.
 \label{eq:body-ex7-cox}
\end{equation}
این مقدار معادل \(\qty{0.0034515}{\farad\per\metre\squared}\) است.

\paragraph*{ضریب اثر بدنه و افزایش آستانه}
از \eqref{eq:threshold-gamma}، ابتدا ضریب بار و سپس \(\gamma\) را
حساب می‌کنیم
\begin{align}
 B_Q&=\sqrt{2(1.6\times10^{-19})(11.7)(8.85\times10^{-14})(2\times10^{17})}
 \notag\\
 &\approx\qty{2.57427e-7}{\coulomb\per\centi\metre\squared\per\volt\tothe{0.5}},\notag\\
 \gamma&=\frac{2.57427\times10^{-7}}{3.4515\times10^{-7}}
 \approx\qty{0.74584}{\volt\tothe{0.5}}.
 \label{eq:body-ex7-gamma}
\end{align}
واحد \(\gamma\) با کنترل ابعاد
\eqref{eq:threshold-gamma-units} سازگار است. جایگذاری در رابطه
کامل \eqref{eq:threshold-body} می‌دهد
\begin{align}
 \Delta V_t&\approx0.74584
 \left(\sqrt{0.85486+4}-\sqrt{0.85486}\right)\,\unit{\volt}\notag\\
 &\approx\qty{0.95377}{\volt},\notag\\
 V_t&=V_{t0}+\Delta V_t\approx0.7+0.95377
 =\qty{1.65377}{\volt}.
 \label{eq:body-ex7-full-result}
\end{align}
رقم‌های میانی برای بازتولید محاسبه حفظ شده‌اند. پاسخ با گردکردن
متناسب با داده‌های تقریبی چنین است
\begin{equation}
 \boxed{\Delta V_t\approx\qty{0.95}{\volt},\qquad
 V_t\approx\qty{1.65}{\volt}}.
 \label{eq:body-ex7-rounded-result}
\end{equation}
این آستانه نسبت به \lr{Source} تعریف می‌شود. برای شروع وارونگی با
\(V_S=\qty{4}{\volt}\)، ولتاژ گیت نسبت به زمین در مدل به
\(V_G=V_S+V_t\approx\qty{5.65}{\volt}\) می‌رسد.
خواسته اصلی تمرین، افزایش \(\qty{0.95}{\volt}\) در خود آستانه است.

\paragraph*{کنترل مستقل و خطای تقریب خطی}
برای کنترل محاسبه، از فرم دقیق
\eqref{eq:body-rationalized-root} استفاده می‌کنیم
\begin{equation}
 \Delta V_t=\frac{\gamma V_{SB}}
 {\sqrt{\phi_s+V_{SB}}+\sqrt{\phi_s}}
 \approx\qty{0.95377}{\volt}.
 \label{eq:body-ex7-independent-check}
\end{equation}
همچنین \(k_\gamma=\gamma/(2\sqrt{\phi_s})\approx0.40334\) است.
تقریب خطی می‌دهد
\(\Delta V_{t,\mathrm{lin}}\approx0.40334(4)
=\qty{1.61335}{\volt}\).
در این بایاس، \(V_{SB}/\phi_s\approx4.68\) و نسبت خطای
\eqref{eq:body-linear-error} برابر
\(R\approx1.6915\) است. پس خطی‌سازی حدود \num{69.2} درصد
بیش‌برآورد می‌کند. مقدار پاسخ از رابطه کامل به دست آمده است.

در جدول \ref{tab:body-linearization-error}، تغییر آستانه با همین
پارامترهای تمرین برای چند بایاس آموزشی مقایسه شده است. با افزایش
بایاس، شیب دقیق کاهش می‌یابد و فاصله تقریب خطی بیشتر می‌شود.
\begin{table}[!htbp]
\centering
\small
\caption{مقایسه مدل کامل و خطی برای پارامترهای تمرین \secnum{۲.۷}. خطای نسبی
از نسبت \((\Delta V_{t,\mathrm{lin}}/\Delta V_t-1)\) به دست آمده است.}
\label{tab:body-linearization-error}
\begin{tabular}{cccc}
\toprule
\(V_{SB}\; (\unit{\volt})\) &
\(\Delta V_t\; (\unit{\volt})\) &
\(\Delta V_{t,\mathrm{lin}}\; (\unit{\volt})\) & خطای نسبی، درصد\\
\midrule
\num{0.05} & \num{0.01988} & \num{0.02017} & \num{1.44}\\
\num{0.10} & \num{0.03922} & \num{0.04033} & \num{2.84}\\
\num{0.50} & \num{0.17855} & \num{0.20167} & \num{12.95}\\
\num{1.00} & \num{0.32619} & \num{0.40334} & \num{23.65}\\
\num{4.00} & \num{0.95377} & \num{1.61335} & \num{69.15}\\
\bottomrule
\end{tabular}
\end{table}

\section*{تمرین \secnum{۲.۸}، تعداد ترانزیستورهای سری در گیت}
\paragraph*{صورت مسئله و فرض مدار}
آیا اثر بدنه تعداد ترانزیستورهایی را که می‌توان در یک گیت \lr{CMOS}
در فرکانس پایین سری کرد محدود می‌کند
\cite[تمرین \secnum{۲.۸}، صفحه ۹۶]{weste}؟
شبکه \lr{Pull-Down} شامل \(N\) ترانزیستور سری \lr{nMOS} را در نظر
می‌گیریم. \(N\) متناهی است، همه گیت‌ها در \(V_{DD}>V_{t0}\) هستند،
بدنه‌ها به زمین وصل‌اند و شبکه \lr{Pull-Up} در این ورودی‌ها خاموش است.
در مدل حالت ماندگار، جریان بار خارجی و نشتی مسیر مقابل را صفر می‌گیریم.

\paragraph*{تخلیه تدریجی گره‌ها و سطح منطقی خروجی}
\lr{Pull-Down} مسیر اتصال خروجی به زمین را فراهم می‌کند.
ترانزیستور پایین سورس متصل به زمین دارد و با گیت بالا روشن است.
این ترانزیستور، بار خازنی گره مجاور را تخلیه می‌کند. با پایین‌آمدن
ولتاژ گره، برای ترانزیستور بعدی \(V_{GS}\) افزایش و \(V_{SB}\)
کاهش می‌یابد. در نتیجه، شرایط هدایت آن تقویت می‌شود و گره بالاتر
نیز تخلیه می‌شود. همین استدلال برای هر تعداد متناهی گره ادامه دارد.

در حد حالت ماندگار بدون جریان خارجی، افت ولتاژ روی مسیر روشن به
صفر میل می‌کند. بنابراین برای سورس ترانزیستور \(i\) داریم
\begin{equation}
 V_{S,i}\longrightarrow0,\qquad
 V_{SB,i}\longrightarrow0,\qquad
 V_{t,i}\longrightarrow V_{t0},\qquad i=1,\ldots,N.
 \label{eq:body-finite-stack-limit}
\end{equation}
طبق \eqref{eq:threshold-body}، اثر بدنه هنگام تخلیه گره‌ها کاهش می‌یابد.
برای این مدل و زمان نشست کافی، اثر بدنه به‌تنهایی حد سختی برای تعداد
ترانزیستورهای سری در رسیدن به صفر منطقی تعیین نمی‌کند.
رفتار دوگان شبکه سری \lr{pMOS} با بدنه در \(V_{DD}\)، رساندن سطح
یک منطقی است. این تفسیر با ویژگی عبور صفر قوی توسط \lr{nMOS} و یک
قوی توسط \lr{pMOS} در خلاصه فصل کتاب سازگار است
\cite[صفحه ۹۵]{weste}.

\paragraph*{زمان نشست و فرض فرکانس پایین}
اثر بدنه در طول گذار، جریان را کاهش می‌دهد. افزایش تعداد قطعات سری
نیز مقاومت مسیر و تعداد گره‌های خازنی را افزایش می‌دهد. بنابراین زمان
لازم برای رسیدن خروجی به محدوده مجاز منطقی بیشتر می‌شود.
برای هر \(N\) انتخاب‌شده، مدت ثابت‌ماندن ورودی‌ها باید برای نشست
همان مدار کافی باشد. پاسخ تمرین درباره حد حالت ماندگار با این فرض
است. در طراحی با فرکانس و تأخیر مشخص، بار خروجی، ابعاد قطعات و
مهلت رسیدن به سطح منطقی، تعداد قابل استفاده را تعیین می‌کنند.
در مدار واقعی، جریان بار و نشتی نیز می‌توانند افت ولتاژ ماندگار ایجاد کنند.

\section*{تمرین \secnum{۲.۹}، علامت بایاس بدنه}
\paragraph*{صورت مسئله و پاسخ با قرارداد پایانه‌ها}
برای افزایش آستانه \lr{nMOS}، باید بایاس بدنه مثبت باشد یا منفی
\cite[تمرین \secnum{۲.۹}، صفحه ۹۶]{weste}؟
در مدل \eqref{eq:threshold-body} با \(\gamma>0\)، افزایش آستانه
از \(V_{SB}>0\) حاصل می‌شود. چون \(V_{SB}=V_S-V_B\)، شرط عمومی چنین است
\begin{equation}
 \boxed{V_B<V_S}.
 \label{eq:body-bias-general-sign}
\end{equation}
اگر سورس به زمین متصل باشد، این شرط به \(V_B<0\) تبدیل می‌شود.
پس پاسخ برای سورس زمین، بایاس منفی بدنه است. این حالت
\lr{Reverse Body Bias (RBB)} نام دارد. با منفی‌ترشدن بدنه نسبت به
سورس، بایاس معکوس پیوند و بار تخلیه افزایش می‌یابند.

برای مثال، اگر \(V_S=\qty{0.5}{\volt}\) و
\(V_B=\qty{0.2}{\volt}\) باشند، \(V_{SB}=\qty{0.3}{\volt}>0\) است.
در این نمونه، بایاس بدنه نسبت به سورس، معکوس است و آستانه افزایش
می‌یابد. علامت بایاس در پاسخ عمومی، نسبت به پتانسیل سورس تعیین می‌شود.

\paragraph*{مثال آموزشی افزایش آستانه با بدنه منفی}
برای سورس زمین و \(V_B=-\qty{0.4}{\volt}\)،
\(V_{SB}=\qty{0.4}{\volt}\) است. اگر
\(V_{t0}=\qty{0.4}{\volt}\)،
\(\gamma=\qty{0.5}{\volt\tothe{0.5}}\) و
\(\phi_s=\qty{0.8}{\volt}\) باشند، داریم
\begin{equation}
 V_t=0.4+0.5\left(\sqrt{0.8+0.4}-\sqrt{0.8}\right)
 \approx\qty{0.5005}{\volt}.
 \label{eq:body-negative-bias-example}
\end{equation}
افزایش آستانه حدود \(\qty{0.1005}{\volt}\) است.
با گیت و سورس ثابت، اضافه‌ولتاژ گیت کاهش می‌یابد و جریان هدایت
در مدل کاهش پیدا می‌کند. اندازه بایاس قابل اعمال در یک فرایند مشخص
باید از اطلاعات همان فرایند تعیین شود.

\section*{کنترل حل مسئله و ادامه جزوه}
برای مسئله‌های اثر آستانه، ابتدا اختلاف ولتاژهای پایانه‌ها و مرجع
\(V_{t0}\) را مشخص کنید. \(N_A\) و \(n_i\) باید واحد یکسان داشته
باشند و محاسبه \(\gamma\) به دستگاه واحد سازگار نیاز دارد.
برای بایاس بدنه، از رابطه کامل رادیکالی شروع کنید. تقریب خطی با
نسبت \(V_{SB}/\phi_s\) و معیار خطای
\eqref{eq:body-linear-error} ارزیابی می‌شود.

پس از محاسبه آستانه مؤثر، \(V_{GS}-V_t\) و مرز ناحیه کاری را
به‌روزرسانی کنید. در مدل‌های جمع‌شده، مرجع بایاس آستانه و سهم
\lr{DIBL} در شیب خروجی را مشخص کنید. در مدار سری، ولتاژ گره میانی
از برابری جریان‌ها به دست می‌آید و با تغییر آستانه جابه‌جا می‌شود.
مسیرهای جریان زیرآستانه، نشتی گیت و نشتی پیوند در بخش
\ref{sec:leakage} همین جزوه آمده‌اند. فصل بعد، تمرین مرتبط با کاهش
آستانه و نشتی را حل می‌کند. اثر عمومی دما و هندسه در بخش‌های
\ref{sec:temperature-dependence} و \ref{sec:geometry-dependence} بررسی شده است.
